%%%%%
%
\section{Rayonnement IR � partir d'un pyrg�om�tre 
(temp\'eratures en unit� quelconque)}

\itemstyle{Op\'erateur~:} {\sf rayon\_IR\_auto\_tempe}
\index{rayon\_IR\_auto\_tempe}

%%%%%OBJET%%%%
\itemstyle{Objet~:}
Calcul du flux de rayonnement infra-rouge \`a partir 
des mesures brutes d'un radiom\`etre (pyrg�om�tre) : tension fournie 
par la thermopile et temp�ratures
du bo�tier et de la coupole.

Cette version d\'etecte automatiquement l'unit\'e dans laquelle
sont exprim�es les temp\'eratures ( C ou K ) et assure
la transformation en Kelvin \'eventuellement n�cessaire.
On suppose que les toutes les temp\'eratures 
pass\'ees en argument ont la m\^eme unit\'e.

%%%%%ARGUMENTS%%%%
\bigskip
%%%%%ENTREES%%%%
\begin{tabarg}
\multicolumn{3}{l}{\titretab {Entr\'ees}}\\
\hline
$V_{th}$ & Vecteur, {\tt float} & Tension thermopile ($V$)\\ %% (ou �V ??)
$T_b$ & Vecteur, {\tt float} & Temp\'erature du bo�tier ($K$ ou $^\circ C$)\\
$T_c$ & Vecteur, {\tt float} & Temp\'erature de la coupole ($K$ ou $^\circ C$)\\
$a$ & Coeff., {\tt float} &  Coefficient de sensiblit\'e de la thermopile 
($Wm^{-2}/V$)\\
$K$ & Coeff., {\tt float} &  Coefficient de correction caract\'eristique de la liaison bo�tier-coupole\\
\end{tabarg}

\medskip
%%%%%SORTIES%%%%
\begin{tabarg}
\multicolumn{3}{l}{\titretab {Sorties}}\\
\hline
$Flx$ & Vecteur, {\tt float} & Flux infra-rouge ($W/m^2$)\\
\end{tabarg}


%%%%%FORMULE%%%%
\itemstyle{Formule}

%%rayon_IR_auto_tempe
\[
Flx =  a.V_{th} + \sigma.T_b^4 + K.\sigma.(T_b^4 - T_c^4)
\]
avec $\sigma$ = constante de Stefan-Boltzmann ($5.6698 10^{-8}$) et $T_b$ et $T_c$ en $Kelvin$.


%%%%%AUTEUR%%%%
\itemstyle{Auteur~:}
CNRM/GMEI/TRAMM

%%%%%REFERENCE%%%%
\itemstyle{R\'ef\'erences~:}
Chepfer-1996 \cite{AIE96-Chepfer}
